\subsection{全体概要：この章で学ぶこと}

この章は、ソフトウェア工学を支えるコンピュータ科学の基礎を扱う。ソフトウェア技術者は、与えられたアルゴリズムを単にコードへ変換するだけではない。要求を読み、構成を設計し、適切なアルゴリズム、通信方式、性能目標、テスト方針、保守方法を選ぶ。そのためには、計算機、データ、プログラム、基本ソフトウェア、データベース、ネットワーク、人間、人工知能の基礎を横断的に理解する必要がある。

到達目標として、読者は次のことを説明できるようになることを目指す。

\begin{itemize}[leftmargin=2em]
\item システム、サブシステム、モジュール、構成要素の関係を説明できる。
\item コンピュータアーキテクチャとコンピュータ組織の違いを説明できる。
\item データ構造、アルゴリズム、計算量の関係を説明できる。
\item プログラミング言語、型、サブルーチン、オブジェクト指向、デバッグの基礎を説明できる。
\item 基本ソフトウェア、データベース、ネットワークがソフトウェア設計へ与える制約を説明できる。
\item 利用者と開発者の人間要因、人工知能・機械学習とソフトウェア工学の関係を説明できる。
\end{itemize}

\par\smallskip
\begin{center}
\centering
\IfFileExists{assets/generated_figures/ch16/fig-ch16-01.png}{\includegraphics[width=0.96\linewidth]{assets/generated_figures/ch16/fig-ch16-01.png}}{\fbox{\parbox[c][48mm][c]{0.9\linewidth}{\centering 画像生成待ち\\16章の全体地図}}}
\par\smallskip
\refstepcounter{figure}\label{fig:ch16-01}図 \thefigure: 16章の全体地図
\end{center}
図 \thefigure は、16章の全体地図を視覚的に整理したものである。
\par\smallskip
